\begin{theorem}[Arbitrarily translated irreducible character pairings]
    \label{thm:peps_translated_character_pairing}
    \lean{Representation.sum_character_inv_mul_character_translate,
        Representation.sum_character_inv_mul_character_translate_eq_zero}
    \leanok
    \uses{thm:peps_isotypic_projectors}
    Let $G$ be a finite group, and let $\sigma,\tau$ be finite-dimensional
    irreducible complex representations. Write $d_\sigma=\dim\sigma$.
    For all $a,b\in G$,
    \[
        \sum_{g\in G}\chi_\sigma((ag)^{-1})\chi_\tau(bg)
        =\begin{cases}
            \dfrac{|G|}{d_\sigma}\chi_\tau(ba^{-1}), & \sigma\cong\tau,     \\
            0,                                       & \sigma\not\cong\tau.
        \end{cases}
    \]
    In particular, if $\chi_\sigma\ne\chi_\tau$, every such translated pairing
    vanishes. This is an auxiliary group identity for
    \cite[charge-detection proof]{Schuch2010PEPS}, lines 2464--2486;
    it makes no assertion about a measurement on the original spins.
\end{theorem}
\begin{proof}\leanok
    Change variables from $g$ to $ag$. Apply character-weighted Schur
    orthogonality to $\sum_g\chi_\sigma(g^{-1})\tau(g)$, multiply by
    $\tau(ba^{-1})$, and take the trace. The sum is $|G|/d_\sigma$ times the
    identity when the irreducible representations are equivalent, and zero
    otherwise. Equality of irreducible characters is equivalent to equivalence
    of their representations.
\end{proof}

\begin{theorem}[Translated charge-vector orthogonality and normalization]
    \label{thm:peps_translated_character_inner_product}
    \lean{Representation.sum_star_character_mul_character_translate,
        Representation.sum_star_character_mul_character_translate_eq_zero,
        Representation.sum_star_character_mul_character_self}
    \leanok
    \uses{thm:peps_translated_character_pairing,
        thm:peps_unitary_character_adjoint}
    Under the preceding hypotheses, suppose $\sigma$ is unitary. Then
    \[
        \sum_{g\in G}\overline{\chi_\sigma(ag)}\chi_\tau(bg)
        =\begin{cases}
            \dfrac{|G|}{d_\sigma}\chi_\tau(ba^{-1}), & \sigma\cong\tau,     \\
            0,                                       & \sigma\not\cong\tau.
        \end{cases}
    \]
    Thus distinct irreducible characters give orthogonal translated vectors,
    independently of $a$ and $b$. Each translated vector of one unitary
    irreducible representation has squared norm
    \[
        \sum_{g\in G}|\chi_\sigma(ag)|^2=|G|.
    \]
    These are the auxiliary character relations used in
    \cite[charge-detection proof]{Schuch2010PEPS}, lines 2464--2486.
    Vectors belonging to the same irreducible representation need not be
    orthogonal for different translations.
\end{theorem}
\begin{proof}\leanok
    Replace the conjugated character by its value at the inverse group
    element. For the norm, put $\tau=\sigma$ and $b=a$; the factor
    $\chi_\sigma(1)=d_\sigma$ cancels the irreducible dimension.
\end{proof}
